\documentclass[11pt]{article}

\usepackage[margin=1in]{geometry}
\usepackage{amsmath,amssymb,amsthm,mathtools}
\usepackage{enumitem}
\usepackage{hyperref}
\usepackage{xcolor}
\usepackage{booktabs}
\usepackage{microtype}

\hypersetup{
  colorlinks=true,
  linkcolor=blue!60!black,
  citecolor=blue!60!black,
  urlcolor=blue!60!black
}

\newtheorem{theorem}{Theorem}
\newtheorem{conjecture}{Conjecture}
\newtheorem{problem}{Open Problem}
\newtheorem{definition}{Definition}
\newtheorem{remark}{Remark}

\newcommand{\Z}{\mathbb Z}
\newcommand{\N}{\mathbb N}
\newcommand{\E}{\mathbb E}
\newcommand{\R}{\mathbb R}
\newcommand{\U}{\mathcal U}
\newcommand{\Cob}{\mathrm{Cob}}
\newcommand{\Col}{\mathrm{Col}}
\newcommand{\ord}{\operatorname{ord}}
\newcommand{\vt}{v_2}
\newcommand{\geom}{\mathrm{geom}}

\title{The Complete Nature of the Collatz Obstruction\\
\large From the Pointwise Valuation Problem to the Shadow Frontier:\\
what was eliminated, what remains, and what new mathematics is required}
\author{John A. Janik\\ \small\texttt{john.janik@gmail.com}}
\date{June 2026}

\begin{document}
\maketitle

\begin{abstract}
This document is the capstone synthesis of a multi-year Collatz program
spanning three campaigns: the pointwise-valuation/SRCE program (the
\texttt{collatz/} manuscripts), the finite-spectral-shadow program
(\texttt{deterministic\_equidistribution\_theory/}), and the Shadow
Frontier (\texttt{shadow\_frontier/}, the endpoint).  No proof of the
conjecture is claimed.  The claim is sharper and, we believe, more
useful: \emph{the obstruction has been located, compressed, priced, and
shown to be invariant under change of framework}.  Three independent
routes --- ours, Tao's almost-all theory, and a contemporaneous
LLM-collaboration program audited herein --- terminate, under explicit
dictionaries, at the same residual statement.  In its most compressed
form the obstruction is one scalar parity sequence: the survival
readings $\vt(3^\ell m-1)=\vt(m-3^{-\ell})$ of the two-clock round
chain, provably i.i.d.\ fair-thirds on Haar fibers, provably
deterministic (and lethal) on ghost-aligned fibers, with every
structured deviation metered by exact clocks, without-replacement
supply, entropy deficits, and an amortized census budget.  What is
missing is exactly one species of theorem: a
\emph{distributional-to-pointwise} (quenched) rigidity statement for a
single self-referential zero-entropy stream.  We trace how each
campaign eliminated its layer of candidate mechanisms, state the
endpoint constraints with their constants
($D^*=1.2463$, $\varepsilon_{\rm bin}=0.449$,
$\varepsilon_{\rm nat}=0.1397$, $R=0.0882$,
$\theta_{12}\le0.97907$), and close with the new mathematics required,
keyed to canonical and contemporary references.
\end{abstract}

\tableofcontents

\section{Accelerated Syracuse notation}

We work with the odd accelerated Syracuse map
\[
  T(x)=\frac{3x+1}{2^{b(x)}},
  \qquad b(x)=\vt(3x+1),
\]
$x$ odd and positive unless stated otherwise.  Along an orbit write
$b_n=b(x_n)$, $B_N=\sum_{n<N}b_n$.  The critical slope is $\log_23$,
and the logarithmic identity
\begin{equation}\label{eq:log-identity}
  \log_2 x_{N}
  =
  \log_2 x_0
  +N\log_2 3-B_N
  +\sum_{n<N}\log_2\Bigl(1+\frac{1}{3x_n}\Bigr)
\end{equation}
is exact.  Mean letter excess $\frac1NB_N>\log_23+\eta$ forces descent
(the C1 mechanism).  For a finite character $\chi'$ the two-point
shadow satisfies the exact keystone identity
\begin{equation}\label{eq:shadow-letter}
  F_{\chi'}(x_n,x_{n+1})=\chi'(x_{n+1})\overline{\chi'(3x_n+1)}
  =\chi'(2)^{-b_n},
\end{equation}
with geometric stationary value
$G(z)=\sum_{b\ge1}2^{-b}z^{-b}$ for $z=\chi'(2)$.
Under Haar measure on $\Z_2$ the letters are i.i.d.\ geometric,
$\Pr[b=k]=2^{-k}$, $\E[b]=2>\log_23$: a typical orbit descends.  The
entire difficulty of the conjecture is the upgrade of this annealed
statement to every single deterministic orbit.

\section{Campaign I: the pointwise valuation problem and SRCE}

The first program (\texttt{collatz/docs/}: the master dependency
document, the pointwise-valuation paper, the stabilized-residue
extinction paper, and the sparse-target pivot note) isolated the
quenched upgrade as a standalone statement and then localized it.

\subsection{The conditional reduction and the annealed theorem}
The honest logical skeleton, recorded in the master dependency
document, is one-directional:
\[
  \mathrm{SRCE}\ +\ \text{Exit Inequality}\ +\ \text{cycle elimination}
  \ \Longrightarrow\ \text{Collatz}.
\]
The \emph{weak Pointwise Valuation Theorem} (PVT) --- the sole
pointwise input --- was proved in its \emph{annealed} form with an
explicit machine-checked Lyapunov function (a correction-weighted,
floor-killed transfer operator of spectral radius
$\le\sqrt3/(2\sqrt2-1)<1$).  The canonical-residue formula
\[
  r_m\equiv 3^{-m}\bigl(2^{D_m}-A_m\bigr)\pmod{2^{D_m+1}}
\]
showed the weak PVT \emph{equivalent} to a deterministic
shrinking-target statement, the \emph{Stabilized Residue Cluster
Extinction} problem (SRCE): no fixed odd $x_0$ admits an infinite
valuation word whose canonical residues stabilize to $x_0$ while the
deficit $D_m-m\log_23$ stays confined.

\subsection{Why Borel--Cantelli fails, and the H23a localization}
The targets have summable density $\asymp3^{-m}$ and the residue map
is per-scale uniform --- yet ordinary Borel--Cantelli is powerless,
for an identified structural reason: the lift digits are nonnegative,
the residue sequence is \emph{monotone}, and hence the deterministic
target hits are \emph{maximally clustered}.  The extremal set was
represented as a peak-rooted inverse-Syracuse tree, modeled as a
killed critical branching random walk in a strip, and SRCE was reduced
to a fixed-modulus Perron--Frobenius gap for the exact deterministic
congruence operator (H23a), with the annealed gap a theorem and the
character-killing mechanisms partially machine-checked in Lean.

\subsection{The lesson of Campaign I}
Every natural mechanism was eliminated by proof or computation:
finite-state and block Lyapunov functions, Baker/convergent
near-cycle control, word-level entropy and complexity,
$\times2\times3$ entropy rigidity, $p$-adic scenery overlap,
restart-density arguments.  The postmortem
(\texttt{why\_collatz\_fails.md}) identified the wall exactly: the
confined counterexample word is \emph{symbolically generic} ---
indistinguishable from a confined renewal word in every word
statistic --- so no annealed or finite-state mechanism can see it.
The sparse-target pivot note drew the correct conclusion: the
remaining obstruction is a \emph{deterministic sparse-target avoidance
problem for a non-invertible affine cocycle}, demanding an
inverse-theorem framework (persistent over-hitting must force finite
algebraic/coboundary structure) rather than another estimate.  This
diagnosis fixed the agenda for everything that followed.

\section{Campaign II: finite spectral shadows}

The second program built the inverse-theorem infrastructure on the
finite (annealed) side and proved the global reduction arrow.

\subsection{Coboundary classification and the finite gap}
The finite coboundary classification $\Cob=\langle\chi_2\rangle$
(parity is the only coboundary direction; the proof uses the
$b\mapsto b+2$ branch comparison and the primitivity of $2$ modulo
powers of $3$) separates genuine shadows from the parity escape.  On
that foundation:

\begin{theorem}[Finite non-coboundary spectral gap, informal]
At every finite conductor level, every primitive non-coboundary
twisted synchronized transfer operator has spectral radius strictly
below the untwisted operator, and the high-conductor sector has a
uniform gap at $s=0$ (numerically
$\sup_c\gamma_c(0)\le(\sqrt3/2)^{1/2}+o(1)\approx0.9306$).
\end{theorem}

The proof required the coboundary classification, conductor collapse,
exact pair-correlation phase identities, a threshold cascade,
bounded-shell flattening, and a root-ball/Hensel analysis of the
degenerate cells.  It is an \emph{annealed} result: the operator mixes
shadows; no single orbit is yet implicated.

\subsection{The keystone arrow and quenched shadow exclusion}
\begin{theorem}[Keystone arrow, informal]
A bad positive orbit (divergent, or trapped in a nontrivial cycle)
must, for some finite non-coboundary $\chi'$, have empirical two-point
shadow averages failing to converge to $G(\chi'(2))$.
\end{theorem}
By \eqref{eq:shadow-letter} the shadows \emph{are} letter characters,
so the statement is finite.  The arrow is one-directional (the trivial
cycle itself deviates), and the corrected global endpoint is
\emph{Quenched Shadow Exclusion}: no non-eventually-trivial positive
orbit carries a persistent non-coboundary shadow --- which is Collatz
in finite-shadow language.

\subsection{Repetition rigidity and the exclusions}
The genuinely new theorem of Campaign II, machine-verified in
Lean~4:

\begin{theorem}[$2$-adic repetition rigidity]
Two positions of a positive orbit carrying the same length-$p$
valuation block, lying within $2^{p+1}$ of each other, are
\emph{equal}; the orbit is then an integer cycle.
\end{theorem}

Combined with the Baker/Rhin two-logarithm bound
($x_{\min}\le Cp^{\tau+1}$, $\tau=13.3$) this kills all cycles in
verified ranges and all low-complexity aperiodic ghosts: Sturmian,
Beatty, morphic, linearly recurrent, polynomial-complexity, and
bounded-discrepancy valuation words are not realizable by divergent
positive orbits.  Mean-visible deviations die by
\eqref{eq:log-identity}.  The $3$-adic-only rigidity hope is
\emph{false} (the all-ones trap is coherent, biased, certificate-free
modulo every $3^r$, and is the ghost $x=-1$); the $2$-adic
realization tower carries the positivity obstruction and is not
auxiliary.

\section{Campaign III: the Shadow Frontier (the endpoint)}

The third campaign (\texttt{shadow\_frontier/}) converted the
surviving-profile description into exact, certified, machine-verified
structure.  Every constant below is computed by exact integer
arithmetic or verified at scale.

\subsection{Exact counting and the phase diagram}
$D$-confinement quantizes the partial sums $B_j$ into integer windows
of size $\le\lfloor2D\rfloor+1$, so confined-word counts $N_D(p)$ are
\emph{exactly} computable (no Perron estimates).  The narrow-band
dichotomy follows: an acyclic orbit dwells at most $CY^{\theta(D)}$
steps in a band $[Y,2^{2D}Y]$, with counting capacity exhausted at
\[
  D^*=1.2463\qquad(\text{band ratio }2^{2D^*}=5.63).
\]
At ratio $\le2$ the dwell bound is \emph{logarithmic} (the dyadic
strip carries exactly one word per phase).  Tilting by the shadow
observable yields the exact $(D,\varepsilon)$ phase diagram, whose
structure is now mapped: the vertical counting wall at $D^*$; the
horizontal \emph{binary bias ceiling}
$\varepsilon_{\rm bin}=\bigl|(2-\log_23)\zeta_3^{-1}+(\log_23-1)\zeta_3-G\bigr|
=0.44903$; the \emph{free-bias basin} below
$\varepsilon_{\rm nat}(D)\to|\E_{\rm cal}-G|=0.1397$ (bias below the
basin floor costs nothing; above it the Legendre tilt $t^*$ is the
exact price, $\partial\theta/\partial\varepsilon=-t^*/\ln2$); genuine
arithmetic kinks at $2D\in\Z$ and Diophantine resonance lines
$D=|m-j\log_23|$; and a first-order \emph{direction-exchange corner}
at $(D,\varepsilon)=(2.11,0.402)$ where the optimal conspiracy
direction jumps from the binary $\{1,2\}$ channel to the
phase-coherent $\{1,4\}$ channel.  A counterexample's banded windows
are bracketed to $[1.2463,\,{\approx}4.08]\times(0,\varepsilon^*(D)]$,
$\varepsilon^*\le0.54$.

\subsection{The two clocks and the universal ghost clock}
Every letter is a $2$-adic depth reading: $b=1$ iff $x\equiv3\ (4)$,
and a maximal $1$-run has length \emph{exactly} $\vt(x+1)-1$ (the
$-1$ clock); $b=2$ iff $x\equiv1\ (8)$, with run length
$\lfloor(\vt(x-1)-1)/2\rfloor$, depth decrement exactly $2$, and a
parity handoff (the $+1$ clock --- the trivial cycle is the $+1$
ghost); $b\ge3$ iff $x\equiv5\ (8)$, reading $b=\vt(x+\tfrac13)$.
This generalizes to \emph{every} signature $\sigma$ of depth $K$:
the phantom root $\rho_\sigma=C_\sigma/(2^K-3^\ell)$ (the catalogued
ghosts $-1,+1,-5,-19/11,-17$ are instances) carries an exact clock ---
alignment drops exactly $K$ per ridden block, rides last exactly
$\lfloor(a-1)/K\rfloor$ blocks, and alignment $a$ at height $x$ costs
$2^a\le|2^K-3^\ell|x+C_\sigma$.  Consequences, all unconditional:
without-replacement depth supply ($2^{k-g}$ pool points per band per
ghost, consumed once by an acyclic orbit); at most $2^{k-3}$ upward
crossings of level $2^k$, ever; a divergent orbit spends
$O(X\log X)$ lifetime steps below $X$.  Amortized over \emph{all}
ghost families of depth $\ge3$, the available drift gain is
$R=0.0882$ bits per letter against the $0.415$-bit budget
($4.7\times$ margin) --- independently recomputed in the audit of the
contemporaneous program of Chang \cite{Chang}.

\subsection{The round chain, the i.i.d.\ coin, and the LTE cascade}
Parsed at $1$-run starts, survival against the letters-$\le2$ regime
is the reading
\[
  v=\vt(3^\ell m-1)=\vt\bigl(m-3^{-\ell}\bigr),
\]
death iff $v$ odd: an alignment reading of fresh seed bits against a
target rotating through $\langle3\rangle\subset\Z_2^\times$.  On Haar
fibers the readings are \emph{exactly} i.i.d.\ (survival $\tfrac13$
per round; the proof composes the orbit-parameterization slope
$2\cdot3^{B}$ with an exact census lemma); apparent finite-modulus
anti-correlations are truncation artifacts.  Surviving rounds drift
\emph{upward} at $+0.6166$ bits per round: the $2$-adic Cantor set of
survivors is measure zero and its positive members, if any, are forced
ascenders.  On ghost-aligned fibers the coin is \emph{deterministic}:
lifting-the-exponent decides Mersenne seeds ($2^g-1$ survives round
one iff $g$ is odd with $\vt(g-1)$ even) and the $3$-smooth families
$2^g3^a\mp1$ in closed form, with second-round readings such as
$\vt\bigl(a'+2-\mathrm{Log}_3(-5)\bigr)$ --- the exponent reading the
$2$-adic logarithm of a ghost --- and every tested seed dies with a
finite certificate.  The binary world is closed unconditionally by the
\emph{entropy deficit}: $\theta_{12}(D)\le H(2-\log_23)=0.97907<1$ at
every band width, so an orbit with letters $\subset\{1,2\}$ has no
phase-diagram refuge anywhere.  The weakest nontrivial case of the
remaining conjecture is precisely Chang's Carry Independence
Conjecture (every $n_0>1$ eventually takes a letter $b\ge3$), which is
the parity-$\equiv1$ extreme of our calibration statement.

\subsection{Faithfulness and the blind spot}
Is this language adequate?  Provably, within its declared scope: the
word coding is the Bernstein--Lagarias conjugacy
\cite{Bernstein,BernsteinLagarias} --- lossless, a homeomorphism onto
the full shift carrying Haar to Bernoulli($2^{-k}$); every
finite-modulus observable of a bounded window factors exactly through
(one residue, the letters); conditional on any word cylinder, residues
at all odd places are exactly equidistributed (only the $2$-adic place
feeds the word); and the empirical hidden-channel test is null
(multiplicative invariants of deep survivors match controls;
celebrated special numbers --- $e^{\pi\sqrt{163}}$'s integer,
worst-case Miller--Rabin composites --- are Collatz-civilians).  The
language's one blind spot is the subset $\N\subset\Z_2$ itself,
imported through four exact interfaces (clocks, counting, the drift
identity, Baker) --- and the blind spot \emph{bites}: a positive seed
has $\log_2x_0$ bits, so after ${\approx}\tfrac23\log_2x_0$ steps
every letter is forced by the $\times3$ carry transducer.  The whole
infinite word has Kolmogorov complexity $\le\log_2x_0+O(1)$ while
rigidity demands near-maximal factor complexity: integer-realizable
words are a \emph{countable, computable} family (of the $2^{30}$
binary words of depth $30$, sixty-six are realized below $10^6$), and
survival champions spend $59$--$76\%$ of their lives in the forced
regime, growing like $2^{D/2.2}$ in their survival depth $D$ and
behaving exactly like best-of-$N$ random search.

\subsection{Convergence of independent frameworks}
Tao's almost-all theory \cite{Tao}, the Chang program \cite{Chang}
(audited piece by piece: its burst--gap calculus \emph{is} the
two-clock parse under an explicit dictionary; its census constant
verifies; its shadow return-time theorem is refuted and replaced by
supply bounds), and this program terminate at the same residual:
almost-all versus all; distributional versus pointwise; annealed
versus quenched.  Three languages, one wall, with explicit
translations --- the language-invariance expected of a faithful
compression of the dynamics.

\section{The complete nature of the obstruction}

The successive localizations compose into one chain:
\[
  \text{Collatz}
  \;\longleftarrow\;
  \text{PVT}
  \;\longleftrightarrow\;
  \text{SRCE}
  \;\longleftarrow\;
  \text{H23a}
  \;\longleftarrow\;
  \text{Quenched Shadow Exclusion}
  \;\longleftarrow\;
  \text{Two-Clock Calibration}
  \;\longleftarrow\;
  \cdots
\]
each arrow an unconditional reduction, terminating in:

\begin{center}
\fbox{\begin{minipage}{0.92\textwidth}
\textbf{The obstruction.}  Show that no positive integer can feed the
round-chain coin even readings with the conspiratorial statistics
required for survival, forever.  Equivalently: the empirical law of
the readings $\vt(m_i-3^{-\ell_i})$ along a single orbit --- a
zero-entropy, computable, self-referential stream, provably
i.i.d.-fair on Haar fibers and provably lethal on every
deterministically analyzable fiber --- cannot deviate from its
distributional law forever.  Every structured deviation channel is
metered: counting ($D^*=1.2463$), bias ceilings
($0.449$, saturation $0.54$), free basin ($0.1397$), alphabet
($\theta_{12}<0.9791$), ghost rides ($\lfloor(a-1)/K\rfloor$ with
height cost $2^a$), supply ($2^{k-g}$, without replacement), census
($R=0.0882$ of $0.415$), and the Kolmogorov budget
($\log_2x_0$ bits, spent in the opening).  The missing theorem is of
distributional-to-pointwise type for one coin.
\end{minipage}}
\end{center}

Three features make this obstruction \emph{complete} as a description:
(i) it is \emph{language-invariant} (three independent frameworks
reach it under explicit dictionaries); (ii) it is \emph{exhaustive
within declared categories} (the finite-shadow language is provably
lossless and complete for congruence mechanisms, with the residual ---
positivity --- explicitly isolated, not compressed away); (iii) it is
of a \emph{known species}: the same distributional-to-pointwise gap
that blocks normality of $\sqrt2$'s digits, effective
Furstenberg $\times2\times3$ rigidity for individual points, and every
``true for almost every $x$, unknown for each $x$'' statement in
metric number theory.  The Collatz wall is not an isolated pathology;
it is the generic wall of quenched arithmetic dynamics, here presented
in its sharpest known finite form.

\section{What new mathematics is required}

The annealed theory is finished; the congruence category is provably
complete; the archimedean imports are exact.  What follows is keyed to
the surviving attack surfaces, with the literature a reader should
enter through.

\subsection{Quenched rigidity for self-referential streams (the coin)}
The central need: an inverse theorem asserting that a computable
zero-entropy stream which persistently deviates from its fiberwise
i.i.d.\ law must exhibit finite algebraic structure (coboundary,
ghost-alignment, or periodicity) --- structure the certificate
machinery then kills.  The nearest existing mathematics:
Furstenberg's $\times2\times3$ disjointness and its measure rigidity
\cite{Furstenberg,Rudolph,Host}, made \emph{effective} in
\cite{BLMV}; Sarnak's M\"obius-disjointness program and the
Bourgain--Sarnak--Ziegler criterion \cite{Sarnak,BSZ} as the model
for ``deterministic sequences cannot correlate with structured
ones''; and the Green--Tao--Ziegler inverse theory \cite{GTZ} as the
template for ``large correlation forces algebraic structure.''  None
applies as stated: the Collatz cocycle is non-invertible, affine, and
self-referential (the stream reads its own output).  The new theorem
must be an \emph{effective, pointwise} rigidity statement for
$2$-adic affine transducers --- a quenched local limit theorem for
one orbit of one machine.

\subsection{Digit rigidity and ``Baker along orbits''}
In the forced regime the letters read binary digits of $3$-power
multiples of the seed, and the deterministic cascades terminate in
readings of $2$-adic linear forms in logarithms
($\vt(N-\mathrm{Log}_3(-5))$ and its deeper analogues).  The entry
points: Stewart's effective lower bounds for digit counts of $3^n$
\cite{Stewart}, Senge--Straus \cite{SengeStraus}, the $p$-adic
linear-forms theory of Yu \cite{Yu}, Bugeaud's monograph
\cite{Bugeaud}, and --- as proof that digit streams \emph{can} be
controlled along sparse deterministic sequences --- Mauduit--Rivat
\cite{MauduitRivat}.  The needed extension is quantitative control of
$\vt$-readings along \emph{orbit-indexed} moving targets: a hybrid of
linear forms in logarithms with the renewal structure of the round
chain (``Baker along orbits''), for which the periodic case
\cite{Baker,Rhin,SimonsDeWeger,Hercher} is the solved boundary case.

\subsection{Banded $S$-unit and subspace rigidity (Target 1)}
High-complexity banded windows supply many relations
$2^{B}x_{n+p}-3^px_n=A_{n,p}$ with $A_{n,p}$ a structured
$\{2,3\}$-unit sum.  The conjecture --- too many coherent relations in
a multiplicative band force the exponent vectors into finitely many
linear patterns, hence repetition or a certificate --- is of
subspace-theorem type: Evertse--Schlickewei--Schmidt \cite{ESS} for
the count of solutions, Corvaja--Zannier methods \cite{CZ} for the
applications template, Amoroso--Viada \cite{AV} for the quantitative
state of the art.  The obstruction to direct application is that the
classical theorems count solutions of \emph{fixed} equations, while
the orbit generates a \emph{growing family} of equations with shared
structure; what is needed is a rigidity theorem for coherent
\emph{families} of $S$-unit relations.

\subsection{Mixed-tower max-plus certificates (Targets 2--3)}
The certificate layer --- Karp maximum cycle means of
$g+\lambda(\phi-\varepsilon)$ over residue traps in $\Z/2^a3^b$, with
the margin $\inf_{\lambda\ge0}(P_H(\lambda)-\lambda\varepsilon)$
implemented and validated in the candidate tools --- needs its
deterministic completion: \emph{every positive-realizable coherent
trap eventually certifies, except the trivial cycle}.  The relevant
mathematics is tropical/max-plus spectral theory \cite{Butkovic} and
ergodic optimization \cite{Jenkinson}; the missing principle is a
projective (lift-coherent) Perron theory for the joint
$(2,3)$-tower, for which the H23a operator comparison of Campaign~I
is the prototype.

\subsection{The automaton route}
The blind-spot theorems reframe the problem as machine behavior:
every integer is a finite seed of one fixed carry transducer, and the
conjecture asserts the machine always halts (reaches the $+1$ ghost).
The relevant boundary results: Conway's undecidability for
generalized Collatz functions \cite{Conway,KurtzSimon} (the family is
Turing-complete; any proof must use the specific arithmetic of
$3x+1$), Cobham-type rigidity and automatic-sequence theory
\cite{Allouche,AdamczewskiBugeaud} (outputs of finite automata cannot
be too structured \emph{and} too generic at once --- precisely the
tension a counterexample must inhabit), and algorithmic randomness
\cite{DowneyHirschfeldt} for the Kolmogorov-versus-factor-complexity
gap.  A genuinely new ``transducer rigidity'' theorem --- the output
stream of a bounded-carry $\times3$ transducer from a finite seed
cannot be a calibrated biased pseudorandom word forever --- would
close the conjecture and would be a fundamental statement about
multiplication itself.

\subsection{Supporting quantitative frame}
Tao's almost-all theorem \cite{Tao} bounds how large the exceptional
set can be; Krasikov--Lagarias density bounds \cite{KrasikovLagarias}
constrain preimage trees; verification stands at $2^{68}$
\cite{Barina} with cycle exclusions through $m\le91$ \cite{Hercher};
Lagarias's surveys and sourcebook \cite{LagariasSurvey,LagariasBook}
remain the canonical maps of the terrain; Wirsching's monograph
\cite{Wirsching} treats the backward tree; Terras \cite{Terras}
founded the density approach.  The contemporary parallel program of
Chang \cite{Chang} is, after our audit, a relabeled fragment of the
same structure with the same wall.

\section{Endpoint}

\begin{center}
\fbox{\begin{minipage}{0.92\textwidth}
There is no longer any reason to search for counterexamples: every
describable candidate class is excluded by an exact meter, the
surviving profile is the output set of one finite machine
conditioned on beating a fair coin forever, and a divergent
counterexample would in any case be computationally unverifiable.
The remaining work is the construction of one theorem ---
distributional-to-pointwise rigidity for the round-chain coin ---
toward which the most promising roads are effective
$\times2\times3$-type rigidity, $p$-adic Baker theory along orbits,
$S$-unit family rigidity, projective max-plus Perron theory, and
transducer rigidity.  The search for the counterexample and the
construction of the proof were the same activity all along: every
emptied hiding place above is now a theorem, and the proof, if it
comes, will be the completed inventory.
\end{minipage}}
\end{center}

\begin{thebibliography}{99}

\bibitem{Baker} A.~Baker, Linear forms in the logarithms of algebraic
numbers, \emph{Mathematika} 13 (1966), 204--216.

\bibitem{Rhin} G.~Rhin, Approximants de Pad\'e et mesures effectives
d'irrationalit\'e, in \emph{S\'eminaire de Th\'eorie des Nombres,
Paris 1985--86}, Progr. Math. 71, Birkh\"auser, 1987, 155--164.

\bibitem{Yu} K.~Yu, $p$-adic logarithmic forms and group varieties
III, \emph{Forum Math.} 19 (2007), 187--280.

\bibitem{Bugeaud} Y.~Bugeaud, \emph{Linear Forms in Logarithms and
Applications}, IRMA Lect. Math. Theor. Phys. 28, EMS, 2018.

\bibitem{ESS} J.-H.~Evertse, H.~P.~Schlickewei, W.~M.~Schmidt, Linear
equations in variables which lie in a multiplicative group,
\emph{Ann. of Math.} 155 (2002), 807--836.

\bibitem{CZ} P.~Corvaja, U.~Zannier, \emph{Applications of Diophantine
Approximation to Integral Points and Transcendence}, Cambridge Tracts
in Math. 212, CUP, 2018.

\bibitem{AV} F.~Amoroso, E.~Viada, Small points on subvarieties of a
torus, \emph{Duke Math. J.} 150 (2009), 407--442.

\bibitem{Furstenberg} H.~Furstenberg, Disjointness in ergodic theory,
minimal sets, and a problem in Diophantine approximation,
\emph{Math. Systems Theory} 1 (1967), 1--49.

\bibitem{Rudolph} D.~Rudolph, $\times2$ and $\times3$ invariant
measures and entropy, \emph{Ergodic Theory Dynam. Systems} 10 (1990),
395--406.

\bibitem{Host} B.~Host, Nombres normaux, entropie, translations,
\emph{Israel J. Math.} 91 (1995), 419--428.

\bibitem{BLMV} J.~Bourgain, E.~Lindenstrauss, P.~Michel,
A.~Venkatesh, Some effective results for $\times a\times b$,
\emph{Ergodic Theory Dynam. Systems} 29 (2009), 1705--1722.

\bibitem{Sarnak} P.~Sarnak, Three lectures on the M\"obius function,
randomness and dynamics, IAS lecture notes, 2011.

\bibitem{BSZ} J.~Bourgain, P.~Sarnak, T.~Ziegler, Disjointness of
M\"obius from horocycle flows, in \emph{From Fourier Analysis and
Number Theory to Radon Transforms and Geometry}, Springer, 2013,
67--83.

\bibitem{GTZ} B.~Green, T.~Tao, T.~Ziegler, An inverse theorem for
the Gowers $U^{s+1}[N]$-norm, \emph{Ann. of Math.} 176 (2012),
1231--1372.

\bibitem{Stewart} C.~L.~Stewart, On the representation of an integer
in two different bases, \emph{J. Reine Angew. Math.} 319 (1980),
63--72.

\bibitem{SengeStraus} H.~G.~Senge, E.~G.~Straus, PV-numbers and sets
of multiplicity, \emph{Period. Math. Hungar.} 3 (1973), 93--100.

\bibitem{MauduitRivat} C.~Mauduit, J.~Rivat, Sur un probl\`eme de
Gelfond: la somme des chiffres des nombres premiers, \emph{Ann. of
Math.} 171 (2010), 1591--1646.

\bibitem{AdamczewskiBugeaud} B.~Adamczewski, Y.~Bugeaud, On the
complexity of algebraic numbers I. Expansions in integer bases,
\emph{Ann. of Math.} 165 (2007), 547--565.

\bibitem{Allouche} J.-P.~Allouche, J.~Shallit, \emph{Automatic
Sequences: Theory, Applications, Generalizations}, CUP, 2003.

\bibitem{Conway} J.~H.~Conway, Unpredictable iterations, in
\emph{Proc. 1972 Number Theory Conf.}, Univ. Colorado, Boulder, 1972,
49--52.

\bibitem{KurtzSimon} S.~A.~Kurtz, J.~Simon, The undecidability of the
generalized Collatz problem, in \emph{TAMC 2007}, LNCS 4484,
Springer, 2007, 542--553.

\bibitem{DowneyHirschfeldt} R.~Downey, D.~Hirschfeldt,
\emph{Algorithmic Randomness and Complexity}, Springer, 2010.

\bibitem{Jenkinson} O.~Jenkinson, Ergodic optimization in dynamical
systems, \emph{Ergodic Theory Dynam. Systems} 39 (2019), 2593--2618.

\bibitem{Butkovic} P.~Butkovi\v{c}, \emph{Max-linear Systems: Theory
and Algorithms}, Springer, 2010.

\bibitem{Tao} T.~Tao, Almost all orbits of the Collatz map attain
almost bounded values, \emph{Forum Math. Pi} 10 (2022), e12.

\bibitem{KrasikovLagarias} I.~Krasikov, J.~C.~Lagarias, Bounds for
the $3x+1$ problem using difference inequalities, \emph{Acta Arith.}
109 (2003), 237--258.

\bibitem{SimonsDeWeger} J.~L.~Simons, B.~M.~M.~de~Weger, Theoretical
and computational bounds for $m$-cycles of the $3n+1$ problem,
\emph{Acta Arith.} 117 (2005), 51--70.

\bibitem{Hercher} C.~Hercher, There are no Collatz $m$-cycles with
$m\le91$, \emph{J. Integer Seq.} 26 (2023), Art. 23.3.5.

\bibitem{Barina} D.~Barina, Convergence verification of the Collatz
problem, \emph{J. Supercomputing} 77 (2021), 2681--2688.

\bibitem{Terras} R.~Terras, A stopping time problem on the positive
integers, \emph{Acta Arith.} 30 (1976), 241--252.

\bibitem{LagariasSurvey} J.~C.~Lagarias, The $3x+1$ problem and its
generalizations, \emph{Amer. Math. Monthly} 92 (1985), 3--23.

\bibitem{LagariasBook} J.~C.~Lagarias (ed.), \emph{The Ultimate
Challenge: The $3x+1$ Problem}, AMS, 2010.

\bibitem{Wirsching} G.~J.~Wirsching, \emph{The Dynamical System
Generated by the $3n+1$ Function}, Lecture Notes in Math. 1681,
Springer, 1998.

\bibitem{Bernstein} D.~J.~Bernstein, A non-iterative $2$-adic
statement of the $3N+1$ conjecture, \emph{Proc. Amer. Math. Soc.} 121
(1994), 405--408.

\bibitem{BernsteinLagarias} D.~J.~Bernstein, J.~C.~Lagarias, The
$3x+1$ conjugacy map, \emph{Canad. J. Math.} 48 (1996), 1154--1169.

\bibitem{Chang} E.~Y.~Chang, Exploring Collatz dynamics with
human-LLM collaboration, arXiv:2603.11066 (2026); audited in the
ghost-clock companion note of this program.

\end{thebibliography}

\end{document}
